\documentclass[11pt]{article}
\usepackage[margin=1in]{geometry}\usepackage{amsmath}\usepackage{hyperref}\usepackage[T1]{fontenc}
\title{Recovering and Auditing a Year of IBM Quantum Experiment Records\\[4pt]\large Thirty jobs, what verifies, and what does not}
\author{Joseph Cyril Dougherty IV\\13th Chamber LLC}
\date{November 16, 2026}
\begin{document}\maketitle
\begin{abstract}We audit every IBM Quantum experiment record recoverable from a year of personal cloud storage: thirty job entries spanning December 2025 to July 2026 across six backends, with register widths from two to forty bits and shot counts from 1,024 to 4,096. Eight job entries survive as original Qiskit Runtime exports; eighteen exist only as transcribed counts in session reports; four could not be retrieved at all. The audit separates what the files establish from what the transcripts assert. Transcribed highlights, including an 86.9 percent Bell-state fidelity and an 80 percent GHZ-4 visibility, remain unverified because the underlying exports were not recovered. The frequently cited CHSH value S = 2.78 cannot be reconstructed: the recovered exports contain computational-basis populations only, with no rotated-basis correlation data and no recovered circuit definitions. The audit also documents four internal inconsistencies in the transcribed record, including shot totals that sum to 1,025 rather than 1,024 and a stated entropy that is incompatible with the stated peak. We argue the package is useful as a worked example of record decay, and we publish the inventory, the discrepancy log and a retention checklist so that future experiments are verifiable by default rather than reconstructable only in principle.\end{abstract}
\section*{1. Scope and method}
The material under audit is a personal archive of IBM Quantum jobs run by the author between January 2025 and October 2026. Records were recovered from Google Drive, OneDrive, exported session transcripts and public web pages. Recovery matched file names, job identifiers, Qiskit Runtime serialisation metadata and count histograms.

Three rules governed the audit. First, a claim is treated as primary evidence only when an original Qiskit Runtime export was recovered and its counts decoded. Second, a claim recorded only in a session transcript is treated as reported, not established. Third, a claim that cannot be traced to any file is dropped entirely, whatever its prominence in earlier writing.

The inventory in Section 3 is the complete set of job entries located. It is a lower bound on what was run and an upper bound on what can be checked today.

\section*{2. What the archive contains}
Thirty job entries were located. Eighteen ran on ibm\_torino, two on ibm\_kingston, and one each on ibm\_marrakesh, ibm\_fez, ibm\_sherbrooke, plus a 133-qubit Torino run and a 156-qubit Heron r2 Fez run recorded under separate labels. Five entries lost their backend metadata during the transfer failures described below.

Twenty-six entries are marked Completed, one Failed, and the remainder carry no usable status. Shot counts cluster at 1,024 (eleven entries) and 4,096 (eleven entries), with two at 2,048 and the rest unknown. Register widths range from two bits to forty bits; the forty-bit registers belong to randomised circuits whose observed distributions are near uniform, with maximum counts of one or two per outcome out of 4,096 shots, which is what uniformly random bitstrings look like at that depth.

Job creation times, where recoverable, run from 15 December 2025 to 6 July 2026. Earlier work from March 2025 survives only as a refetch log recording RuntimeJobNotFound on both instances, which is itself a useful datum: job records on hosted quantum services expire.

\section*{3. Primary evidence: the eight recoverable exports}
Eight entries were recovered as original exports with decodable counts. Four carry forty-bit registers on ibm\_torino and ibm\_marrakesh from 16 February 2026, with maximum outcome counts of two. One ibm\_kingston four-bit job from 27 March 2026 shows a peak of 151 counts at outcome 0011, that is 14.7 percent over sixteen states, with the remaining fifteen states between 88 and 107 counts. One Kingston job with split registers records maximum counts of 88, 57 and 281 across a five-bit, five-bit and two-bit register. One Torino job survives as a digitised histogram with a dominant outcome near 445 counts, about 43.5 percent. One Fez job failed at submission and produced metadata only.

Two decoding cautions apply throughout. Byte order is ambiguous in the transcribed summaries: the same Kingston peak is called 0011 under most-significant-bit-first decoding and 1100 under the opposite convention. And a failed job is not an experiment, so the Fez entry carries no result.

\section*{4. Reported but unverified results}
Eighteen entries exist only as transcribed counts. Among the most cited: a Bell-state run on ibm\_torino reporting 480 counts at 00, 410 at 11, 60 at 01 and 75 at 10, described as 86.9 percent fidelity; a GHZ-4 run reporting roughly 450 counts each at 0000 and 1111, described as 80 percent visibility; a Quantum Volume 8-qubit run across 256 states; a five-qubit Navier-Stokes table with the all-ones outcome at 0.81 percent; and a nine-qubit run with 859 counts at the all-zero outcome, or 20.97 percent.

These are internally plausible for the stated hardware and shot counts, and several are consistent with known depolarising behaviour. They are nonetheless reported values: without the original exports the audit cannot confirm them, and it does not treat them as established.

\section*{5. The CHSH claim does not reconstruct}
Earlier writing cited a CHSH parameter of S = 2.78 with an uncertainty near 0.04, attributed to two specific jobs. The recovered exports cannot support that number. A CHSH parameter requires coincidence counts in several rotated measurement bases with known classical bit mappings; what was recovered are computational-basis populations from circuits whose definitions were not retained. No set of recovered files determines S.

The record is also inconsistent about it. Documented values across the year include 2.78, 2.82, 2.875, 2.41, 1.87 and 0.875. One of those, 1.87, sits below the classical bound of 2.0 and was reported as such by the author at the time. A quantity that ranges across those values in the same archive is a quantity that was never pinned to a single measurement protocol.

For the same reason, the widely repeated claim of a 348-fold error suppression is not supported by the recovered exports. It may well have been observed; the archive does not let anyone check.

\section*{6. Four discrepancies in the transcribed record}
First, the Bell counts sum to 1,025 against a stated 1,024 shots. Second, a three-qubit run reports 700 counts from 1,024 at a single outcome, 68 percent, while also reporting 2.958 bits of entropy over eight states; those two statements cannot both hold. Third, the register decoding direction is stated inconsistently between summaries, which changes which bitstring is called the peak. Fourth, three job identifiers quoted in one transcript belong to IBM's public documentation tutorial and are not the author's jobs.

None of these invalidate the work. All of them matter, because each is the kind of error that survives into a paper, a slide deck or a funding application and is then very expensive to retract.

\section*{7. Why the record decayed}
Four mechanisms account for almost all of the loss. Hosted services expire job records, so a job identifier is not a durable citation. Exports were stored in personal cloud folders in at least four locations, with duplicate file names that differ in content. Results were transcribed into prose in session reports, and the prose outlived the JSON. And circuit definitions were kept separately from results, so even where counts survived, the measurement bases needed to interpret them did not.

\section*{8. A retention checklist}
\begin{itemize}\item Hash every result export at the moment of retrieval and record the digest beside the job identifier.\end{itemize}
\begin{itemize}\item Store the circuit that produced each result with the result, not in a separate folder.\end{itemize}
\begin{itemize}\item Keep the measurement basis and the classical bit mapping for every run; without them an entangled-state result is uninterpretable.\end{itemize}
\begin{itemize}\item Write the shot count once, derive every percentage from it, and state the arithmetic check in the note.\end{itemize}
\begin{itemize}\item Mirror the whole set to an immutable public archive at submission time, so a claim can be checked after the account expires.\end{itemize}
\begin{itemize}\item Separate files from transcripts in the folder structure and never quote a transcript as though it were a file.\end{itemize}
\section*{9. Conclusion}
A year of quantum experiments produced thirty job entries, eight verifiable exports and one published inventory. The audit's value is not that the numbers are large; it is that the boundary between verified and reported is now written down. The Bell and GHZ transcriptions are consistent with the hardware and worth re-running properly. The CHSH figure should be treated as withdrawn until a multi-basis measurement produces it from recoverable data. Anyone can check all of this: the inventory, the discrepancy log and the derived counts ship with this paper.

\section*{10. Expanded recovery, November 2026}

A second recovery pass was run after first publication, targeting sources the original audit could not reach: the Notion workspace, a Perplexity session export, Grok session exports, and the wider Drive tree. The pass indexed 293 Drive objects matching quantum, IBM, Qiskit, QSAM or session keywords, against the 30-entry inventory taken earlier. That number is objects, not experiments: it includes duplicates, notes, generated reports and archives.

A job-context file of 105 KB recovered from Drive contains 188 job-identifier mentions resolving to 25 distinct identifiers. Twenty-four were already in the inventory. One is new: d5098l7g0u6s73dbio30, a chemistry run on ibm_torino associated with an aspirin and COX-bin Hamiltonian. It has no recovered counts, so it enters the inventory as reported-only, which changes the composition to thirty-one entries: eight with original exports, nineteen reported-only, four unrecoverable.

Three identifiers that appear in session text, d2n1d2cg59ks73c6vu3g, d2n18dehb60s73cvalg0 and d2ff2cms6qcs738f4lb0, belong to IBM's public documentation tutorial and are excluded. They were excluded before; the second pass confirms that exclusion rather than extending it, and the audit now states the rule explicitly: an identifier is evidence only when it can be tied to a submission made by the author.

Notion was searched for the first time. It holds 95 pages and 5 databases. Fifty-five pages discuss quantum work, largely architecture notes for QAI2, the Ghost Fingerprint platform and Q-Consciousness. The one database that sounds relevant, Experiment Results, is an unrelated 2022 template of landing-page tests, so Notion contributes narrative context rather than measurements. No job counts were recoverable from it.

The Perplexity export is a two-entry archive with no extractable text, so it adds nothing numeric. The position of Section 5 is unchanged by the second pass: the recovered computational-basis populations still cannot produce a CHSH parameter, and S = 2.78 remains withdrawal-worthy rather than confirmed. What the second pass changes is completeness, not conclusions, which is exactly what a re-audit should be able to say.

\section*{11. Physical recordings of the runs, November 2026}

The audit previously recorded that circuit definitions were not retained. That finding is now qualified. A Drive sweep for video files returned eleven distinct phone recordings, duplicates excluded, totalling roughly 1.9 GB and ranging from 29 to 178 seconds, whose file names carry capture dates of 14 and 15 December 2025. Those dates coincide with the Turin job batch that produced the only dense forty-bit exports in the archive.

Frames extracted from the recordings were read directly. They show the IBM Quantum Platform composer: circuit setup, the QPU selection step naming a processor and its qubit count, a pending-jobs queue, and the confirmation line that a job has been submitted. The archive therefore contains contemporaneous screen evidence of submission, not only after-the-fact transcripts. What it still does not contain is the exported counts for those same jobs.

That distinction is the whole point of this paper, so it is worth stating precisely. A recording proves an event happened at a time, on a named backend, with a circuit on screen. It does not supply the outcome distribution, and it cannot be re-analysed, re-binned or re-plotted. For the audit's purposes, recordings move a claim from unsupported to corroborated, not from reported to verified.

A manifest ships with this version: file name, byte size, duration, Drive identifier and MD5 checksum for all eleven recordings, so a reader can confirm they are examining the same files. A parallel set of album photographs from the same sessions was retrieved and inspected.

Two practical consequences follow. First, future experiments should capture the exported counts beside the phone recording, in the same folder, with the same checksum routine, because one without the other is either unverifiable or unanalysable. Second, a recording is the cheapest possible insurance against record decay: a phone in front of a screen for two minutes outlived every JSON file in this archive except eight.

\section*{12. Contemporaneous screenshots, November 2026}

A second photographic set was recovered from a shared album: seventeen screenshots of the IBM Quantum console captured during the December 2025 sessions. Their value is specificity. They show the configure-settings panel with region Washington DC (us-east) and a shot count of 1024, the QPU selection step listing the Heron family, the workloads table filtered to that region with a status column, and a results view rendering computational basis states with measured frequencies for a six-qubit circuit.

One frame is decisive for chronology. The usage panel records the job as created at 18:50 on 14 December 2025, pending for zero seconds, and entering progress at the same minute. Another shows the service instance, the resource group and a job identifier prefix matching a job already in this inventory. The submission is therefore anchored in time and in service, by an artifact created at the moment of submission rather than reconstructed afterwards.

The screenshots do not close the gap identified in Section 5. A rendered frequency chart is a picture of a histogram, not a count export: the numbers visible are not the full outcome distribution, the digitisation is manual, and the circuit and basis definitions remain outside the archive. What the set does establish is that the runs were real, on the named backend, at the recorded minute, with the stated shot count.

This suggests the immediate action for the author while the service still holds the records: re-export the counts for the job identifiers in this inventory from the instance shown in the screenshots, then checksum them beside the manifest. Records on hosted services expire, and a re-export now converts several reported-only entries into verified ones at no experimental cost. A manifest of the seventeen screenshots, with checksums and a description of each frame, ships with this version.

\section*{13. Recovered circuit definition and backend selection, November 2026}

Section 5 recorded that circuit definitions were not retained, which is why no recovered file could be re-interpreted. A further set of seven screenshots qualifies that statement. One frame shows the composer's OpenQASM 2.0 panel with the header and register declarations legible: OPENQASM 2.0, the standard include, a four-qubit quantum register and a four-bit classical register. The same frame carries the working notes of the session that built it: start from a Bell state, apply a Hadamard to the first qubit, apply a controlled-NOT from the first to the second, then measure.

This is partial but consequential recovery. Two facts follow that were previously unknown. The measurement basis is computational and the classical mapping is one-to-one with the quantum register, so any count table belonging to this circuit is directly interpretable. And the intended state is a Bell pair embedded in a four-qubit register, which means the four-bit outcomes split into two groups whose populations carry the entanglement signature. A reported-only entry with four-bit counts can now be checked against a known circuit rather than a guess.

What remains missing is the gate list after the pair, since the panel is photographed mid-build rather than on completion. The four-qubit GHZ visible in another entry still has no recovered definition.

The same album contains an observation that explains the shape of the whole archive. The backend selection dialog lists two processors with their live queues: ibm_marrakesh at 156 qubits with 14,038 pending jobs, and ibm_torino at 133 qubits with zero pending jobs. The inventory is dominated by Torino. That is not a design preference, it is a queue artefact: the fastest available backend was chosen at the moment of submission, repeatedly.

Stated plainly, this is a documented selection effect and it belongs in the limitations of any comparative claim. An archive that samples one processor because it had an empty queue cannot support cross-backend reproducibility, and it cannot support a claim that a result was hardware independent. Recording the queue state at submission, as these screenshots accidentally did, is the fix, and it is worth adding to the retention checklist in Section 8.

\section*{14. Outcomes and identifiers visible on screen, November 2026}

A further album of twenty-two console screenshots arrived after the previous revision, held back briefly by the album service: the link returned an empty album twice before the images propagated. That is itself a small lesson for the retention checklist, and it is the reason the audit does not treat an empty shared folder as evidence of loss.

Two frames in this set change what can be said without qualification. The first is a results readout stating that the three most frequent outcomes of a 1024-shot run account for 345 counts between them, that is 33.7 percent of all shots. This is the first numeric outcome statement in the entire recovery that was recovered from an image rather than from prose. The adjacent frequency panel lists sixteen states, 0000 through 1111, which fixes the register at four qubits and confirms the computational basis already inferred from the recovered OpenQASM panel in Section 13.

The second is a frame showing the service instance reference together with a job identifier that matches an entry already in this inventory, the 14 December 2025 submission recorded there as job B. That identifier was previously attributed on the strength of a transcript. It is now corroborated by an on-screen console reference created at the time, which is a stronger class of evidence than a later retelling.

The limits are worth stating as carefully as the gains. A photographed aggregate does not name the three outcomes it counts, so it cannot be re-binned, cannot be compared against a model fit, and cannot be checked against a second run by anyone. Sixteen bars are visible but only one carries a legible value, and a bar without a legend cannot be tied to a specific state. The strongest honest reading is: for one 1024-shot four-qubit run in December 2025, roughly a third of the shots fell into three outcomes, which is consistent with a noisy Bell pair and inconsistent with a uniform distribution.

The remaining frames are session artefacts rather than measurements: an AI answer describing seven reviewed sources, a spaces page titled Quantum Software Breakthroughs, and Perplexity and browser chrome. They are valuable for reconstructing what the author was doing and when, and they belong in the manifest. They are not experimental evidence, and this audit does not promote them to that status.

\section*{Data availability}
The job inventory (30 entries), the evidence report (all nine sections, including the coverage gaps) and the derived counts ship with this paper as CSV and Markdown, together with the LaTeX source of this document. Nothing in the package requires an IBM account to read.

\section*{References}\begin{enumerate}
\item IBM Quantum, Save jobs and results, quantum.cloud.ibm.com/docs/en/guides/save-jobs (tutorial job identifiers).
\item Qiskit Runtime primitives and result serialisation, docs.quantum.ibm.com.
\item Clauser, Horne, Shimony and Holt, Proposed experiment to test local hidden-variable theories, Phys. Rev. Lett. 23, 880 (1969).
\item Greenberger, Horne and Zeilinger, Bell theorem without inequalities, in Bell's Theorem, Quantum Theory and Conceptions of the Universe (1989).
\end{enumerate}\end{document}